\documentclass[conference]{IEEEtran}

\usepackage[T1]{fontenc}
\usepackage[utf8]{inputenc}
\usepackage{amsmath,amssymb}
\usepackage{booktabs}
\usepackage{array}
\usepackage{tabularx}
\usepackage{microtype}
\usepackage{xurl}
\usepackage[hidelinks]{hyperref}
\usepackage{listings}

\urlstyle{same}
\raggedbottom
\setlength{\emergencystretch}{1em}

\newcolumntype{Y}{>{\raggedright\arraybackslash}X}

\lstset{
  basicstyle=\ttfamily\scriptsize,
  columns=fullflexible,
  keepspaces=true,
  breaklines=true,
  breakatwhitespace=false,
  showstringspaces=false,
  frame=single,
  rulecolor=\color{black},
  aboveskip=6pt,
  belowskip=6pt
}
% The default listing frame does not require a color package.
\lstset{rulecolor={}}

\newcommand{\repo}{
  \url{https://github.com/tyreefranklinjr/
  Mission-Computer--Communications-Platform}
}
\newcommand{\revision}{
  \nolinkurl{e04681fdf181b30d5230861b9ff7cf7fef1f30c1}
}

\title{A Simulation-Based Mission Computer and Resilient
Communications Platform:\\
Architecture, Fault Injection, and Evaluation}

\author{
  \IEEEauthorblockN{Tyree Franklin Jr.}
  \IEEEauthorblockA{
    Computer Science Student\\
    Western Governors University\\
    \url{https://github.com/tyreefranklinjr}
  }
}

\begin{document}
\maketitle

\begin{abstract}
This paper examines a mission-computer simulation and a separate
resilient communications simulator implemented in C++20, C11,
and Rust. The mission computer combines synthetic inertial and
satellite-navigation measurements, a four-state Kalman estimator,
grid-based A* planning, proportional-integral-derivative command
generation, and watchdog-driven state management. The
communications simulator implements sequenced messages,
acknowledgement timeouts, bounded transmission attempts,
duplicate suppression, and score-based link failover.
Evaluation combines source inspection with numerical analysis
of committed artifacts at a fixed repository revision.
All 500 recorded fault scenarios satisfy the implemented
final-state checks. Across 1,500 trajectory samples, planar
position root-mean-square error is 0.138 m. Twenty communications
trials record simulated failover latencies of 22--41 ms, with
a mean of 39 ms. Source inspection also identifies limits on
these findings: mission-state safing does not inhibit actuator
commands, controller output does not drive the simulated plant,
and random-number scaling changes the intended communications
loss model. The study documents the implemented behavior,
quantifies the archived results, and defines verification
criteria for integrated and hardware-based evaluation.
\end{abstract}

\begin{IEEEkeywords}
Mission computers, software-in-the-loop simulation, fault
injection, Kalman filtering, A* search, watchdogs,
communications resilience.
\end{IEEEkeywords}

\section{Introduction}

Mission-oriented software coordinates estimation, planning,
command generation, and health supervision. Its response to
a fault depends on both the local algorithm and the interfaces
between these functions. A useful prototype must make those
interfaces inspectable and associate its claims with specific
evidence.

The Mission-Computer--Communications-Platform repository
implements these concerns in two executables~\cite{repo}.
The C++ mission computer represents four logical nodes within
one process. A separate Rust executable examines message
delivery during a primary-link outage.

This study asks three questions:
\begin{enumerate}
  \item Do the committed fault scenarios produce the mission
  states required by the implemented checks?
  \item What estimation errors and link-switching latencies
  appear in the committed artifacts?
  \item Which implementation and experimental limits prevent
  stronger system-level conclusions?
\end{enumerate}

The contribution is an engineering case study rather than a
new estimation or search algorithm. It connects implementation
details to measured outputs and identifies the additional
evidence needed to evaluate integrated operation.

\section{Scope and Evidence}

\subsection{Revision and Review Method}

The evaluated revision is:
\begin{center}
\begin{minipage}{0.97\columnwidth}
\footnotesize\revision
\end{minipage}
\end{center}

The review covers the mission-computer source, Rust
communications source, two C++ test programs, build
configuration, trial script, and four committed CSV
files~\cite{mission,comms,data,tests}.

The evidence has three distinct forms. Source inspection
establishes what is implemented. Calculations over committed
CSV files establish what those artifacts record. Proposed
changes and acceptance criteria describe work that has not
been performed in this study.

The executables were not rebuilt or rerun for this manuscript.
The numerical findings are therefore a reanalysis of archived
outputs, not an independent reproduction of execution.

\subsection{Execution Boundary}

Both executables run as ordinary software processes.
The implementation does not use physical sensors, a CAN
transceiver, a radio, or an ESP32-C3. FreeRTOS is not linked.

The mission model advances in 20 ms logical steps; the
communications model advances in 1 ms logical steps.
These increments define simulated time. They do not measure
operating-system scheduling accuracy or hardware deadlines.

The executables do not exchange operational messages.
Their outputs support separate subsystem observations,
not an integrated communications-and-control result.

\section{System Architecture}

\subsection{Mission-Computer Execution}

Every mission tick invokes the sensor, navigation, control,
and supervisor functions in that order. Table~\ref{tab:nodes}
summarizes their responsibilities.

\begin{table}[t]
\caption{Logical Mission-Computer Nodes}
\label{tab:nodes}
\centering
\footnotesize
\begin{tabularx}{\columnwidth}{@{}lYY@{}}
\toprule
Node & Responsibility & Output \\
\midrule
Sensor &
Scripted motion and noisy measurements &
IMU/GNSS frames and heartbeat \\
Navigation &
State estimation and route planning &
Position estimate and heartbeat \\
Control &
Waypoint-directed command calculation &
Actuator frame and heartbeat \\
Supervisor &
Health checks and state transitions &
Mission state \\
\bottomrule
\end{tabularx}
\end{table}

Sequential execution simplifies inspection and scenario
generation. It does not represent preemption, concurrent
tasks, independent clocks, or a scheduler that remains
responsive when a real task blocks.

The injected system-overrun fault skips control-node
execution. It models missing work and heartbeats rather
than excessive processor consumption.

\subsection{C/C++ Messaging}

The C library implements a 64-entry ring buffer of CAN-style
frames. Each frame contains an identifier, priority field,
payload length, and eight-byte payload. Publication and
reception take constant time. Storage is fixed, and rejected
publications increment a drop counter.

Reception follows FIFO order. The priority field does not
implement CAN arbitration. Serialization time, physical
errors, bus arbitration, and bus-off recovery are outside
the model.

Floating-point values are packed using memory copies.
A deployment across heterogeneous processors would need
an explicit byte order and numeric representation.

\subsection{Mission-State Management}

The state machine contains five states:
\texttt{INIT}, \texttt{CALIBRATING},
\texttt{AUTONOMOUS}, \texttt{DEGRADED},
and \texttt{SAFE\_STOP}.

Initialization enters calibration. Completion of the
calibration interval permits autonomous operation.
A communications outage moves autonomous operation into
the degraded state; restored communications permits
recovery. An unhealthy core-system indication selects
\texttt{SAFE\_STOP}, which remains terminal for that run.

This state selection does not gate the control function.
The implementation can continue generating commands after
entering \texttt{SAFE\_STOP}. The demonstrated property is
therefore terminal state selection, not actuation inhibition.

\section{Algorithmic Design}

\subsection{State Estimation}

The estimator maintains position and velocity:
\begin{equation}
\mathbf{x}_k =
\begin{bmatrix}
x_k & y_k & v_{x,k} & v_{y,k}
\end{bmatrix}^{T}.
\end{equation}

For acceleration input
$\mathbf{u}_k=[a_x,a_y]^T$, prediction is
\begin{equation}
\hat{\mathbf{x}}^-_k =
F\hat{\mathbf{x}}_{k-1}+B\mathbf{u}_k,
\end{equation}
with
\begin{equation}
F =
\begin{bmatrix}
1&0&\Delta t&0\\
0&1&0&\Delta t\\
0&0&1&0\\
0&0&0&1
\end{bmatrix},
\quad
B =
\begin{bmatrix}
\frac{1}{2}\Delta t^2&0\\
0&\frac{1}{2}\Delta t^2\\
\Delta t&0\\
0&\Delta t
\end{bmatrix}.
\end{equation}

GNSS measures position:
\begin{equation}
H =
\begin{bmatrix}
1&0&0&0\\
0&1&0&0
\end{bmatrix}.
\end{equation}

The covariance prediction, gain, and state correction are
\begin{align}
P^-_k &= FP_{k-1}F^T+Q,\\
K_k &= P^-_kH^T(HP^-_kH^T+R)^{-1},\\
\hat{\mathbf{x}}_k &=
\hat{\mathbf{x}}^-_k+
K_k(\mathbf{z}_k-H\hat{\mathbf{x}}^-_k).
\end{align}

The implementation uses $Q=0.02I_4$, $R=0.5I_2$, and
\begin{equation}
P_k=(I-K_kH)P^-_k.
\end{equation}

Although the source class is named \texttt{Ekf}, its
transition and measurement models are linear.
A four-state linear Kalman estimator is the precise
description of the implemented mathematics.

Fixed-size matrices give constant work per update for this
state dimension. Further evaluation should check covariance
symmetry, positive semidefiniteness, and behavior under
ill-conditioned measurement updates.

The simulated GNSS noise is uniform with amplitude 0.5 m.
Its ideal per-axis variance is
$0.5^2/3\approx0.0833\ \mathrm{m}^2$.
The filter instead uses a measurement-covariance entry of
0.5. The tuning is therefore not a direct match to the
simulator's noise variance.

\subsection{Grid-Based Planning}

The planner uses four-connected, unit-cost moves and
\begin{equation}
f(n)=g(n)+h(n),
\end{equation}
where
\begin{equation}
h(n)=
\sqrt{(x_n-x_g)^2+(y_n-y_g)^2}.
\end{equation}

Euclidean distance is a consistent lower bound under these
movement rules. The search follows the A* formulation
introduced by Hart, Nilsson, and Raphael~\cite{astar}.

The mission instance uses a $20\times20$ grid with a 2 m
coordinate increment. Replanning is requested every
250 ticks, or 5 s of simulated time, and when no path
has been established.

The open list is a vector. Selecting and removing its
lowest-score entry requires linear work in the frontier
size. On this bounded-degree grid, worst-case search
cost can grow quadratically with the number of cells.

The occupancy array holds 4,096 cells, but constructor
arguments are not checked against that capacity.
A reusable version should reject invalid dimensions
before initializing the array.

\subsection{Command Generation}

The controller applies proportional, integral, and
derivative terms with output saturation:
\begin{equation}
u_k =
\operatorname{clip}
\left(K_pe_k+K_iI_k+K_dD_k,-2,2\right).
\end{equation}

The gains are $K_p=1.2$, $K_i=0.02$, and $K_d=0.15$.
The integral is clamped to $[-4,4]$.

The controller input is
\begin{equation}
e_k =
\operatorname{atan2}
(y_{\mathrm{wp}}-\hat y_k,\,
 x_{\mathrm{wp}}-\hat x_k).
\end{equation}

This is a waypoint bearing. The code does not subtract
vehicle heading, so the input is not a conventional
heading-tracking error. Controller output also does not
alter the scripted trajectory. The current evidence
supports command calculation, not closed-loop tracking
or stability.

\subsection{Communications}

The Rust simulator uses a priority heap, an in-flight
message map, a receiver sequence set, and a rolling
link-success history~\cite{comms}.

New messages are generated every 10 simulated milliseconds.
Successful attempts have a 5 ms acknowledgement delay;
unsuccessful attempts time out after 20 ms. Each message
receives at most five total transmission attempts.
The constant named \texttt{MAX\_RETRIES} includes the
initial attempt in this limit.

The primary link is forced down over
$[5000,35000)$ ms. For $m\leq6$ resolved outcomes,
its score is
\begin{equation}
S=\frac{1}{m}\sum_{i=1}^{m}d_i,
\qquad d_i\in\{0,1\}.
\end{equation}

Failover requires at least three recorded outcomes and
a primary score below 0.5.

The heap sends more urgent messages first and uses
older sequence numbers to break ties. However, the
simulator drains the entire outgoing queue in one
iteration. Without a service-capacity constraint,
this ordering does not establish priority-dependent
latency under congestion.

\section{Evaluation and Results}

\subsection{Fault Outcomes}

Each recorded scenario represents 500 mission ticks,
or 10 simulated seconds, with one injected fault.
Table~\ref{tab:faults} reports the recalculated outcomes.

\begin{table}[t]
\caption{Committed Fault-Sweep Results}
\label{tab:faults}
\centering
\footnotesize
\begin{tabularx}{\columnwidth}{@{}Yrrl@{}}
\toprule
Category & Cases & Pass & Stop delay \\
\midrule
Network outage & 125 & 125 & N/A \\
Skipped control task & 125 & 125 & 100 ms \\
CAN publication fault & 125 & 125 & 100 ms \\
Short GNSS dropout & 64 & 64 & N/A \\
Long GNSS dropout & 61 & 61 & 800 ms \\
\midrule
Total & 500 & 500 & \\
\bottomrule
\end{tabularx}

\smallskip
\parbox{\columnwidth}{\scriptsize
Delays use simulated time and are measured from the
injection tick. N/A denotes a case that should not
enter the terminal stop state.}
\end{table}

All 500 final-state checks pass~\cite{data}.
Of these, 311 require \texttt{SAFE\_STOP}; the other
189 require avoiding that state. This is not a result
of 500 successful safe stops.

For tolerated faults, the existing oracle accepts
either \texttt{AUTONOMOUS} or \texttt{DEGRADED} as the
final state. It does not separately require evidence
of both degradation and recovery.

Heartbeat health fails when its age exceeds five ticks.
The recorded five-tick reaction interval starts at the
first omitted heartbeat. The last successful heartbeat
occurred one tick earlier, so heartbeat age is six
ticks at detection.

GNSS safing occurs on the forty-first consecutive
missing sample. The recorded interval from the first
missing-sample tick is 40 ticks.

The cases establish conformance within the generated
scenarios. They do not estimate operational reliability
or cover simultaneous faults.

\subsection{Estimation Accuracy}

For each committed trajectory sample, planar error is
\begin{equation}
e_k =
\sqrt{(\hat x_k-x_k)^2+(\hat y_k-y_k)^2}.
\end{equation}

The planar root-mean-square error is
\begin{equation}
\mathrm{RMSE} =
\sqrt{\frac{1}{N}\sum_{k=1}^{N}e_k^2}.
\end{equation}

Table~\ref{tab:position} uses all 1,500 samples,
including initialization.

\begin{table}[t]
\caption{Position Error in the Committed Trajectory}
\label{tab:position}
\centering
\footnotesize
\begin{tabularx}{\columnwidth}{@{}Yr@{}}
\toprule
Metric & Value \\
\midrule
Planar RMSE & 0.1379 m \\
Mean planar error & 0.1224 m \\
95th percentile, nearest-rank & 0.2412 m \\
Maximum planar error & 0.3621 m \\
Samples below 0.3 m & 1,490 / 1,500 \\
Fraction below 0.3 m & 99.33\% \\
\bottomrule
\end{tabularx}
\end{table}

CSV positions are rounded to four decimal places.
The measurements describe one scripted trajectory and
noise realization. They do not establish performance
under multipath, inertial bias, asynchronous sampling,
or real sensor outages.

\subsection{Link Switching}

Switching latency is
\begin{equation}
L_{\mathrm{switch}} =
t_{\mathrm{switch}}-t_{\mathrm{outage}}.
\end{equation}

Table~\ref{tab:latency} summarizes the 20 committed trials.

\begin{table}[t]
\caption{Simulated Link-Switching Latency}
\label{tab:latency}
\centering
\footnotesize
\begin{tabularx}{\columnwidth}{@{}Yr@{}}
\toprule
Metric & Value \\
\midrule
Trials & 20 \\
Minimum & 22 ms \\
Mean & 39 ms \\
Median & 41 ms \\
Maximum & 41 ms \\
Trials below 250 ms & 20 / 20 \\
\bottomrule
\end{tabularx}
\end{table}

These are virtual-time measurements. Switching the
selected link is distinct from restoring successful
application-level delivery.

The seed-7 event log records 4,394 transmission attempts,
4,000 unique-delivery events labeled \texttt{acked},
393 timeouts, and 140 suppressed duplicates.
Its single failover occurs at 5,041 ms.
No failback event appears.

A new message is transmitted at the terminal timestamp,
40,000 ms. Its normal acknowledgement would occur after
execution ends. Delivery accounting must distinguish
such unresolved messages from losses.

\section{Threats to Validity}

\subsection{Probability Model}

The Rust generator extracts \texttt{state >> 33},
retaining 31 bits, but divides by the maximum 32-bit
unsigned value. Its output therefore occupies
approximately $[0,0.5)$ rather than $[0,1)$.

Under a uniform-state approximation, configured loss
thresholds of 0.03 and 0.05 correspond to approximately
6\% and 10\% event probabilities. The late-arrival
threshold of 0.20 corresponds to approximately 40\%.
Priority proportions also change.

The recorded results remain observations of the
committed implementation. They do not validate its
intended probability configuration. Correcting the
generator requires regenerating the communications
baseline.

\subsection{Repeatability and Recovery}

A fixed seed alone does not ensure identical event
traces. Iteration over randomized \texttt{HashMap}
storage can change the order in which resolved outcomes
enter the rolling history. Deterministic replay should
sort due events before applying them.

Failback requires the primary score to exceed 0.9.
The inactive primary receives no explicit health probes,
so restoring its availability does not itself provide
new successful outcomes. Recovery requires an explicit
probing policy and tests.

\subsection{Safety and Integration}

Terminal mission-state selection does not establish
a safe physical effect. Actuator commands are not
inhibited, the plant does not respond to those commands,
and the Rust simulator does not transport mission
traffic.

The current artifacts cannot establish closed-loop
control performance, integrated command continuity,
or physical stopping behavior.

\subsection{Test Coverage}

The existing tests cover estimator convergence and
small-grid planning~\cite{tests}. They do not establish
complete state-transition coverage, malformed-frame
handling, queue saturation behavior, or bounded
communications memory use.

The planner test checks an upper bound on path length
and verifies obstacle avoidance. It does not explicitly
check adjacency between consecutive cells. A stronger
oracle should verify each transition and compare path
cost with an independent shortest-path reference.

\section{Verification Roadmap}

The next evaluation should address correctness before
expanding the operating envelope.

\subsection{Immediate Gates}

\begin{enumerate}
  \item \textbf{Actuation inhibition:}
  Define a maximum delay from fault detection to command
  inhibition. Check that nonzero commands remain
  inhibited until an explicit reset.

  \item \textbf{Probability correction:}
  Verify the generator's range and event frequencies.
  Regenerate every communications artifact after
  correcting the scaling.

  \item \textbf{Stronger fault checks:}
  Verify the complete transition sequence, reaction
  deadline, final state, and actuator behavior for
  each scenario.
\end{enumerate}

\subsection{Integration Gates}

\begin{enumerate}
  \item \textbf{Closed-loop plant:}
  Feed commands into an explicit plant model, include
  vehicle heading in the tracking error, and measure
  route error and stopping behavior.

  \item \textbf{Operational interface:}
  Exchange sequenced mission messages between the
  executables. Measure loss, duplication, stale
  delivery, and application-level outage duration.

  \item \textbf{Recovery and resource bounds:}
  Test primary-link probing, repeated outage cycles,
  queue saturation, outstanding-message limits, and
  duplicate-history expiration.
\end{enumerate}

\subsection{Target Evaluation}

Hardware testing should record execution times, jitter,
deadline misses, and watchdog behavior under load.
Each result should identify its commit, configuration,
toolchain, target, and input seeds.

A comparative communications experiment should run
fixed-primary transmission, score-based failover, and
probe-assisted recovery under identical offered loads
and outage schedules. Endpoints should include unique
delivery ratio, application-level outage duration,
latency percentiles, unresolved messages, and peak
queue occupancy.

\section{Reproducibility}

The following commands describe the repository workflow.
They were inspected but not executed for this manuscript.

\begin{lstlisting}
git clone https://github.com/tyreefranklinjr/Mission-Computer--Communications-Platform.git
cd Mission-Computer--Communications-Platform
git checkout e04681fdf181b30d5230861b9ff7cf7fef1f30c1

cmake -S . -B build \
  -DCMAKE_BUILD_TYPE=Release
cmake --build build
ctest --test-dir build --output-on-failure

mkdir -p evaluation
./build/mission_computer sim \
  evaluation/trajectory.csv
./build/mission_computer fault-sweep \
  500 evaluation/faults.csv

cd comms_link
cargo build --release --locked
bash run_failover_trials.sh
\end{lstlisting}

The Rust trial script writes into \texttt{samples/}
and updates the seed-7 event log. Reproduction should
use a disposable checkout to preserve the archived
baseline.

An execution record should retain compiler versions,
operating system, build flags, command outputs,
raw CSV files, and analysis procedures.

\section{Conclusion}

The platform implements estimation, grid search,
message exchange, mission-state supervision, and
link-failover simulation in C++, C, and Rust.
Its committed artifacts contain 500 satisfied
final-state checks, a planar position RMSE of
0.138 m on one trajectory, and simulated switching
latencies of 22--41 ms across 20 trials.

These results have a defined scope. The current
implementation selects a stop state without inhibiting
actuation, calculates commands without closing the
plant loop, and evaluates communications separately
from mission traffic. The probability-model defect
also requires a new communications baseline.

The repository provides an inspectable foundation for
further experiments. Connecting safety decisions to
actuation, correcting the probability model, and
strengthening the test oracles are the next steps
toward a defensible integrated evaluation.

\begin{thebibliography}{00}

\bibitem{repo}
T. Franklin Jr.,
``Mission-Computer--Communications-Platform,''
GitHub repository, revision
\revision, accessed Oct. 20, 2026.
[Online]. Available:
\url{https://github.com/tyreefranklinjr/Mission-Computer--Communications-Platform/tree/e04681fdf181b30d5230861b9ff7cf7fef1f30c1}

\bibitem{mission}
T. Franklin Jr.,
``Mission-computer implementation,''
\texttt{src/}, repository revision specified in [1].
Files include \texttt{main.cpp}, \texttt{ekf.hpp},
\texttt{astar.hpp}, \texttt{pid.hpp},
\texttt{watchdog.hpp}, \texttt{mission\_state.hpp},
\texttt{fault\_injector.hpp},
\texttt{imu\_gnss\_sim.hpp}, and
\texttt{can\_bus.c/.h}.

\bibitem{comms}
T. Franklin Jr.,
``Communications simulator,''
\texttt{comms\_link/src/main.rs},
\texttt{comms\_link/src/link.rs}, and
\texttt{comms\_link/run\_failover\_trials.sh},
repository revision specified in [1].

\bibitem{data}
T. Franklin Jr.,
``Committed simulation artifacts,''
\texttt{samples/}, repository revision specified in [1].
Files:
\texttt{fault\_sweep\_results.csv},
\texttt{ekf\_trajectory.csv},
\texttt{comms\_link\_failover\_latencies.csv}, and
\texttt{comms\_link\_log.csv}.

\bibitem{tests}
T. Franklin Jr.,
``Estimator and planner tests,''
\texttt{tests/test\_ekf.cpp} and
\texttt{tests/test\_astar.cpp},
repository revision specified in [1].

\bibitem{astar}
P. E. Hart, N. J. Nilsson, and B. Raphael,
``A formal basis for the heuristic determination
of minimum cost paths,''
\emph{IEEE Transactions on Systems Science and Cybernetics},
vol. 4, no. 2, pp. 100--107, Jul. 1968,
doi: \href{https://doi.org/10.1109/TSSC.1968.300136}
{10.1109/TSSC.1968.300136}.

\end{thebibliography}

\end{document}
